\chapter{Experiment 2: Semantic Conflict in a Master Prompt (RQ3)}
\label{chap:experiment-conflict}

A master prompt can become a shared project artifact that is passed to a later
session or contributor. In this setting, a new instruction may conflict with a
requirement already settled in the artifact.

Experiment~2 uses the settled master prompts of both projects to examine this
problem. It compares three conflict-handling conditions, including a separate
\emph{spec-guided} governance instruction written for the project-handoff
setting. This instruction follows the same general conflict-handling principle
as the RQ2 method, but it is a different prompt.

\section{Defining a Semantic Conflict}
\label{sec:conflict-setting}

A \emph{semantic conflict} occurs when two requirements cannot both hold. For
example, a settled master prompt may state:

\begin{quote}
\emph{Invoices should have a due date 15 days after they are created.}
\end{quote}

A later prompt may then request:

\begin{quote}
\emph{Add a reminder that gets sent 3 days before an invoice's 30-day due
date.}
\end{quote}

The later prompt does not directly ask to change the due date, but it assumes
a 30-day value that conflicts with the settled 15-day requirement. This
differs from a clear revision such as:

\begin{quote}
\emph{Change the invoice due date to 30 days.}
\end{quote}

A clear revision explicitly replaces the earlier value. A semantic conflict
instead occurs when a later prompt assumes or introduces an incompatible value
without clearly stating that the settled requirement should be changed.

The model cannot reliably determine whether the user intended a revision or
simply overlooked the earlier requirement \citep{lahiri2026intent}. Silently
selecting either value may therefore change the user's intent
\citep{liu2025condorcet}.

Semantic conflicts may be difficult to notice because the requirements can use
different wording and appear in separate parts of the master prompt. Unlike
line-based tools such as \texttt{git}, which detect overlapping text changes,
they require comparing requirements by meaning
\citep{pugachev2025codecrdt}.

\section{Experimental Design}
\label{sec:conflict-data}

Experiment~2 uses the hand-built ground-truth master prompts of both projects
as settled project artifacts. Every run follows the same sequence: the agent
receives a fixed handoff prompt, the settled master prompt \(M\) for the
selected project, and then the later prompts one at a time.

\subsection{Handoff and Starting Master Prompts}

After the pre-experiment setup, the same neutral handoff prompt is used in
every condition:

\begin{quote}
\emph{Here is the consolidated result of the earlier prompts for this project.
Continue the project from here. I will then send more prompts, one at a time.}
\end{quote}

The corresponding ground-truth master prompt is then provided as \(M\). The
hard-project \(M\) represents the settled result of its 45-prompt history,
while the simple-project \(M\) represents the settled result of its 12-prompt
history.

Using the ground truths ensures that all agents begin with the same correct
project state. The experiment can therefore focus on conflict handling rather
than errors introduced during earlier consolidation.

\subsection{Later Prompts}

The hard-project case contains four conflict prompts and one compatible
positive control. The simple-project case contains one conflict prompt and one
compatible positive control. All prompts are sent one at a time.

The conflicts are written as ordinary feature requests that assume something
incompatible with \(M\), without clearly asking to replace the settled
requirement. The five conflicts are summarized in
\cref{tab:conflict-taxonomy}, and their exact wording is reproduced in
\cref{app:conflict-prompts}.

\begin{table}[htbp]
  \centering
  \caption{The five semantic conflicts used in Experiment~2.}
  \label{tab:conflict-taxonomy}
  \small
  \begin{tabular}{@{}llp{4.5cm}p{4.5cm}@{}}
    \toprule
    Project & Type & Settled in \(M\) & Assumed by the later prompt \\
    \midrule
    Hard & Value & The invoice due date is 15 days. & A reminder is sent three days before an invoice's 30-day due date. \\
    Hard & Scope & Only vendors use the recycle bin; customers are permanently deleted. & The customers page has a recycle-bin tab and a Restore button. \\
    Hard & Layout & The application uses a top navigation bar. & A button collapses and expands a left sidebar. \\
    Hard & Permission & Staff cannot delete products. & An audit entry records when a staff member deletes their own product. \\
    Simple & Value & The navigation contains four buttons. & All five navigation buttons wrap onto two rows on small screens. \\
    \bottomrule
  \end{tabular}
\end{table}

In the simple-project case, \(M\) defines four navigation buttons:
Dashboard, Policies, History, and Contact. The later prompt refers to five
buttons without requesting an additional button, creating a conflict between
the settled and assumed counts.

The print-button prompt in the hard project and the footer prompt in the
simple project serve as positive controls. Neither conflicts with its
corresponding master prompt and both should be added normally.

\Cref{fig:conflict-example} illustrates the hard-project due-date conflict
under the three conditions.

\begin{figure}[htbp]
  \centering
  \begin{tikzpicture}[
      box/.style={draw, rounded corners, align=left, inner sep=5pt,
        font=\small},
      settled/.style={box, fill=gray!12, text width=11.4cm},
      newp/.style={box, fill=blue!8, text width=11.4cm},
      cond/.style={box, fill=white, text width=4.3cm,
        font=\footnotesize},
      bad/.style={cond, fill=red!10},
      mid/.style={cond, fill=yellow!15},
      good/.style={cond, fill=green!10},
      lbl/.style={font=\bfseries\small},
      condlbl/.style={font=\bfseries\footnotesize},
      arr/.style={-{Latex[length=2.2mm]}, thick},
    ]

    \node[settled] (m) {
      \lbl{Settled in \(M\)}\\[2pt]
      ``Invoices should have a due date 15 days after they are created.''
    };

    \node[newp, below=0.55cm of m] (p) {
      \lbl{Later instruction}\\[2pt]
      ``Add a reminder that gets sent 3 days before an invoice's 30-day due
      date.''\\
      \emph{The prompt requests a feature but assumes a different due date.}
    };

    \node[bad, below=1.05cm of p, xshift=-4.7cm] (c1) {
      \condlbl{Consolidate-then-ask}\\[2pt]
      The settled value may be replaced before the agent is asked whether a
      conflict exists.
    };

    \node[mid, below=1.05cm of p] (c2) {
      \condlbl{Do-not-resolve}\\[2pt]
      The conflicting values are kept separate for the user to decide.
    };

    \node[good, below=1.05cm of p, xshift=4.7cm] (c3) {
      \condlbl{Spec-guided}\\[2pt]
      The settled value remains in the master prompt and the conflicting
      assumption is recorded.
    };

    \draw[arr] (m) -- (p);
    \draw[arr] (p.south) -- (c1.north);
    \draw[arr] (p.south) -- (c2.north);
    \draw[arr] (p.south) -- (c3.north);
  \end{tikzpicture}

  \caption{The due-date conflict under the three experimental conditions.}
  \label{fig:conflict-example}
\end{figure}

\section{Conflict-Handling Conditions}
\label{sec:conflict-conditions}

Claude, Codex, and Gemini are tested under the three conditions shown in
\cref{tab:conflict-conditions}. For each project, every condition receives the
same setup prompt, handoff prompt, settled master prompt \(M\), and later
prompts. The conditions differ only in how conflicts are handled.

\begin{table}[htbp]
  \centering
  \caption{The three conflict-handling conditions.}
  \label{tab:conflict-conditions}
  \begin{tabular}{@{}lp{9cm}@{}}
    \toprule
    Condition & Procedure \\
    \midrule
    \emph{Consolidate-then-ask} & The agent consolidates \(M\) and the later prompts using the \emph{strict} instruction. It is asked about conflicts only after the new master prompt has been produced. \\
    \emph{Do-not-resolve} & The \emph{strict} instruction is extended with a rule that conflicting values must not be selected or merged. The conflict is listed separately for the user to decide. \\
    \emph{Spec-guided} & The separate Experiment~2 governance instruction is active while the later prompts are processed. Each prompt is checked against the settled master prompt when it arrives. \\
    \bottomrule
  \end{tabular}
\end{table}

The comparison examines whether checking for conflicts only after
consolidation is sufficient, whether a direct no-resolution rule prevents
silent replacement, and whether continuous checking preserves settled
requirements as later prompts arrive.

\section{The Spec-Guided Conflict-Handling Method}
\label{sec:conflict-detector}

The \emph{spec-guided} condition uses the final Experiment~2 governance
instruction reproduced in \cref{lst:conflict-governance}. Unlike the RQ2
governance instruction, which maintains \texttt{prompts.md} and
\texttt{spec.md} from the beginning of a project, this instruction begins
with an already-settled master prompt handed to a new session.

Each later prompt is checked against \(M\). Compatible requirements are added,
while clear revisions update the affected requirement. If a later prompt
cannot coexist with a settled requirement and does not clearly request a
change, the settled requirement remains in the master prompt. Both conflicting
requirements are listed under \texttt{\#\# Unresolved conflicts}, and the user
is asked which one should be kept.

When consolidation is requested, the agent returns the conflict-free
requirements under \texttt{\#\# Master prompt} and any unresolved conflicts
in a separate section. The decision is based on the meaning of the
requirements rather than fixed words such as ``change'' or ``instead.''

Semantic Commit provides a related approach in which new instructions are
checked against requirements retrieved from an intent store
\citep{vaithilingam2025semanticcommit}. Retrieval is not required here because
the complete master prompt fits within the model's context. The governance
instruction therefore checks the settled artifact directly.

\section{Evaluation Criteria}
\label{sec:conflict-evaluation}

Each conflict is manually reviewed using the criteria in
\cref{tab:conflict-evaluation}. The review checks whether the agent detects the
conflict, avoids resolving it automatically, preserves the settled
requirement, introduces no new details, and identifies both source
requirements.

\begin{table}[htbp]
  \centering
  \caption{Manual evaluation criteria for each semantic conflict.}
  \label{tab:conflict-evaluation}
  \small
  \begin{tabular}{@{}lp{9.2cm}@{}}
    \toprule
    Criterion & Expected behavior \\
    \midrule
    Conflict detection & The incompatible requirements are reported as a conflict. \\
    No automatic resolution & The agent does not select or merge the conflicting values without user input. \\
    Preservation & The settled requirement remains in the master prompt. \\
    No invention & No new value, condition, or technical detail is introduced. \\
    Source identification & The settled requirement and the conflicting later prompt are identified. \\
    \bottomrule
  \end{tabular}
\end{table}

The print-button and footer prompts are evaluated separately as positive
controls. They should be added without being reported as conflicts.

The experiment is a qualitative demonstration. The criteria define how the
outputs are inspected, but they are not used to calculate a general
conflict-detection score.

\section{Threats to Validity}
\label{sec:conflict-threats}

Experiment~2 is a small qualitative study based on five constructed conflicts
and two compatible controls. Although the cases cover value, scope, layout,
and permission conflicts, they do not represent every inconsistency that may
occur in real projects. Their direct wording may also make them easier to
detect than longer or more ambiguous conflicts.

Model outputs are non-deterministic, and only a small number of runs were
performed. The observed behavior may therefore vary across repeated sessions.

Finally, the conflicts, expected behavior, and outputs were reviewed by the
same author who designed the experiment. Without an independent rater,
interpretation and evaluation bias cannot be ruled out.
